% !TeX program = pdflatex
% SEC-PD — Sprechzettel / Präsentationsgrundlage
% Nicht automatisch kompiliert. Später z. B.:
%   pdflatex -output-directory=docs docs/secpd_sprechzettel.tex
\documentclass[11pt,a4paper]{article}

\usepackage[T1]{fontenc}
\usepackage[utf8]{inputenc}
\usepackage[ngerman]{babel}
\usepackage{geometry}
\usepackage{booktabs}
\usepackage{longtable}
\usepackage{tabularx}
\usepackage{array}
\usepackage{xcolor}
\usepackage{hyperref}
\usepackage{enumitem}
\usepackage{listings}
\usepackage{tikz}
\usepackage{fancyhdr}
\usepackage{titlesec}
\usepackage{microtype}
\usepackage{csquotes}
\usepackage{parskip}

\geometry{margin=2.2cm}
\hypersetup{colorlinks=true,linkcolor=blue!50!black,urlcolor=blue!50!black}
\usetikzlibrary{arrows.meta,positioning,fit,calc}

\definecolor{secpdblue}{HTML}{1B4F72}
\definecolor{secpdgray}{HTML}{F4F6F7}
\definecolor{sprechbg}{HTML}{EAF2F8}
\definecolor{warnbg}{HTML}{FDEBD0}
\definecolor{codebg}{HTML}{F7F9F9}

\pagestyle{fancy}
\fancyhf{}
\fancyhead[L]{\small\textcolor{secpdblue}{SEC-PD \textperiodcentered\ Sprechzettel}}
\fancyhead[R]{\small Credit Risk Innovation}
\fancyfoot[C]{\thepage}
\renewcommand{\headrulewidth}{0.4pt}

\titleformat{\section}{\Large\bfseries\color{secpdblue}}{\thesection}{0.8em}{}
\titleformat{\subsection}{\large\bfseries\color{secpdblue!80!black}}{\thesubsection}{0.7em}{}

\newcommand{\code}[1]{\texttt{\detokenize{#1}}}
\newcommand{\file}[1]{\texttt{\small #1}}
\lstset{
  basicstyle=\ttfamily\small,
  backgroundcolor=\color{codebg},
  frame=single,
  rulecolor=\color{black!15},
  breaklines=true,
  columns=fullflexible,
  keepspaces=true,
}

\newenvironment{sprech}{%
  \par\medskip\noindent
  \begin{tikzpicture}
    \node[fill=sprechbg, draw=secpdblue!40, rounded corners=2pt,
          inner sep=7pt, text width=\dimexpr\textwidth-16pt\relax]
      {\textbf{\textcolor{secpdblue}{Sprechen.}}
}{%
      };
  \end{tikzpicture}
  \par\medskip
}

\newenvironment{achtung}{%
  \par\medskip\noindent
  \begin{tikzpicture}
    \node[fill=warnbg, draw=orange!60!black, rounded corners=2pt,
          inner sep=7pt, text width=\dimexpr\textwidth-16pt\relax]
      {\textbf{Achtung.}
}{%
      };
  \end{tikzpicture}
  \par\medskip
}

\title{\textbf{SEC-PD}\\[0.3em]
\large Ausfallscore aus 10-K MD\&A, Finanzkennzahlen und 8-K-Events\\[0.6em]
\large Sprechzettel: Architektur, Features, Code-Karte}
\author{Praktikumsbeitrag \textperiodcentered\ Credit Risk Innovation}
\date{Stand: Branch \texttt{feat/distress-text-features}, August 2026}

\begin{document}
\maketitle
\tableofcontents
\newpage

% =============================================================================
\section{30-Sekunden-Pitch}
% =============================================================================

\begin{sprech}
SEC-PD schätzt je Firm-Year eine \textbf{12- bzw.\ 36-Monats-Ausfallwahrscheinlichkeit}
aus drei Quellen: XBRL-Finanzkennzahlen, Distress-Sprache in der 10-K MD\&A und
8-K-Frühindikatoren. Das Insolvenz-8-K (Item~1.03) ist \textbf{nur Label}, nie Feature.
Das Modell ist ein kalibrierter Random Forest, offline-fähig (Bank-Gateway) und
über versionierte \texttt{.joblib}-Bundles übertragbar.
\end{sprech}

\textbf{Was es ist:} eigenständiges Python-Modul, eine Woche Praktikum, unabhängig
vom Zielsystem der Bank.\\
\textbf{Was es nicht ist:} keine regulatorische PD, kein Marktpreis-Modell, kein
vollständiges Credit-Rating. Chapter~11 ist ein Default-\emph{Proxy}.

\textbf{Leitmetrik:} PR-AUC und Top-10\,\%-Capture, nicht ROC allein --- Defaults
sind selten.\\
\textbf{Produktionslabel (aktuell):} 36-Monats-Insolvenz, ohne Delisting, ohne
vor-2004-Item-3 (unzuverlässige Submissions-Metadaten).

% =============================================================================
\section{Architektur}
% =============================================================================

\subsection{Datenfluss (Home $\rightarrow$ Training)}

\begin{center}
\begin{tikzpicture}[
  font=\small,
  box/.style={draw=secpdblue, rounded corners=2pt, align=center,
              inner sep=5pt, fill=white, text width=3.4cm},
  arr/.style={-{Latex}, secpdblue, thick},
]
\node[box] (zen) {Zenodo 17121948\\[-1pt]{\scriptsize MD\&A + AAER}};
\node[box, right=0.55cm of zen] (conv) {\texttt{convert\_zenodo}\\[-1pt]{\scriptsize Firm-Year-Tabelle}};
\node[box, right=0.55cm of conv] (edg) {EDGAR\\[-1pt]{\scriptsize companyfacts + 8-K}};
\node[box, below=0.85cm of conv] (tr) {\texttt{train.py}\\[-1pt]{\scriptsize Features $\rightarrow$ Split $\rightarrow$ RF}};
\node[box, below=0.85cm of zen] (fin) {\file{features/financial.py}\\[-1pt]{\scriptsize Ratios, Trends, SIC}};
\node[box, below=0.85cm of edg] (txt) {\file{features/textual.py}\\[-1pt]{\scriptsize LLM + Keywords}};
\node[box, below=0.85cm of tr] (out) {\texttt{models/*.joblib}\\[-1pt]{\scriptsize + Benchmarks}};

\draw[arr] (zen) -- (conv);
\draw[arr] (conv) -- (edg);
\draw[arr] (conv) -- (tr);
\draw[arr] (edg) -- (tr);
\draw[arr] (fin) -- (tr);
\draw[arr] (txt) -- (tr);
\draw[arr] (tr) -- (out);
\end{tikzpicture}
\end{center}

Drei Quellen, ein Join-Schlüssel \texttt{(cik, fyear)}:

\begin{enumerate}[leftmargin=1.4em]
  \item \textbf{Zenodo} liefert 10-K-MD\&A-Texte (und optional AAER-Fraud-Fenster).
        Keine Finanzkennzahlen.
  \item \textbf{EDGAR companyfacts} liefert XBRL-Positionen $\rightarrow$ kanonisches
        Finanz-Panel.
  \item \textbf{EDGAR submissions} liefert 8-K-Metadaten: Insolvenz-\emph{Label} und
        Event-\emph{Features}, strikt getrennt.
\end{enumerate}

\subsection{Zwei Welten: Home und Bank}

\begin{tabularx}{\textwidth}{@{}l X X@{}}
\toprule
 & \textbf{Home} (Internet) & \textbf{Bank} (air-gapped) \\
\midrule
Daten holen & Zenodo + EDGAR-APIs & nur lokale CSVs / Cache \\
LLM & OpenAI \texttt{gpt-5.6-luna} oder Mock & internes Gateway (\texttt{SECPD\_LLM\_MODE=bank}) \\
Replay & Cache unter \file{data/cache/llm/} & dieselben JSON-Dateien, keine API \\
Install & \texttt{pip install -r requirements.txt} & \file{scripts/install\_offline.sh} + Wheels \\
Python & 3.12 getestet & 3.14: Pins ab Commit \texttt{c54e756} \\
\bottomrule
\end{tabularx}

\begin{sprech}
Der LLM-Cache ist der Trick für die Bank: einmal scoren, JSON committen, auf der
Box nur noch Replay. Training und Scoring brauchen dann kein LLM mehr.
\end{sprech}

\subsection{Modell-Modi}

\begin{itemize}[leftmargin=1.2em]
  \item \textbf{financial} --- nur \texttt{fin\_*} + \texttt{evt\_*} (Baseline).
  \item \textbf{combined} (Produktion) --- Financial plus Text in \emph{einem}
        Random Forest (Option~A).
  \item \textbf{ensemble} --- separates Text-Logit-Modell, Aggregation im Logit-Raum
        (Option~B, Gewichte typisch 0{,}6 Finanz / 0{,}4 Text).
\end{itemize}

Estimator: \texttt{RandomForestClassifier} (\texttt{class\_weight=balanced\_subsample},
300 Bäume), davor Median-Imputation. Optional \texttt{CalibratedClassifierCV}
(sigmoid bei wenigen Positiven, isotonic wenn $\ge 100$ Positive; Group-Splits über CIK).

\begin{achtung}
Alte \texttt{.joblib} aus sklearn~1.4.2 nicht auf sklearn~1.7.2 laden --- nach dem
3.14-Pin-Bump neu trainieren. Bundles stempeln Python- und sklearn-Version.
\end{achtung}

% =============================================================================
\section{Label-Logik (was \emph{nicht} ins Modell darf)}
% =============================================================================

\subsection{Default-Definition}

\[
\texttt{label\_default}=1
\quad\Longleftrightarrow\quad
\text{Insolvenz-8-K liegt in }(\text{reporting date},\;\text{reporting date}+H].
\]

$H$ ist 12, 24 oder 36 Monate. Produktionswahl: \textbf{$H=36$, nur Bankruptcy}.

Zusätzlich:
\begin{itemize}[leftmargin=1.2em]
  \item Firm-Years \textbf{nach} der Insolvenz werden gedroppt (kein Post-Event-Lernen).
  \item Rechtszensierte Zeilen (Horizont ragt über das Beobachtungsende) werden gedroppt.
  \item Vor dem Regimewechsel 2004-08-23: Item~3 der Submissions-API ist unzuverlässig
        (Walmart-ASDA-Amendment wurde fälschlich als Insolvenz gezählt).
        Default: \texttt{trust\_legacy\_regime=False}.
  \item Training nur ab GJ~\textbf{2009} (\texttt{require\_financials}): davor keine
        flächigen XBRL-Companyfacts.
\end{itemize}

\subsection{Leakage-Regeln (prüfbar, in Tests abgesichert)}

\begin{tabularx}{\textwidth}{@{}l X@{}}
\toprule
Regel & Begründung \\
\midrule
Item 1.03 / alt-3 nie als Feature & Das \emph{ist} das Label. \\
Delisting nicht als Feature, wenn es Label ist & \texttt{feature\_exclusions\_for\_labels} \\
Item 4.02 Restatement standardmäßig aus & zu nah am Fraud-Label \\
Event-Fenster endet am 10-K-Filing & sonst Look-ahead \\
Text-Features vor dem Split, ohne Fitting & deterministisch, daher kein Leak \\
Temporal-Split, plus Rolling-Origin + Group-Split & keine Zukunft, keine Firm-Memorization \\
\bottomrule
\end{tabularx}

\begin{sprech}
Wenn jemand ``Daten-Snooping'' fragt: das härteste Leakage-Risiko wäre, die
Insolvenzmeldung selbst als Feature zu nutzen. Das ist explizit verboten und
per Test abgesichert. Delisting als \emph{Label} haben wir geprüft --- höhere
Basisrate, aber schlechteres Ranking; deshalb bleibt reine Insolvenz.
\end{sprech}

% =============================================================================
\section{Feature-Katalog}
% =============================================================================

Drei Präfixe: \texttt{fin\_} Finanzen, \texttt{evt\_} 8-K-Zähler, \texttt{txt\_}/\texttt{llm\_} Text.
Combined = alle drei. Financial = \texttt{fin\_}+\texttt{evt\_}.

\subsection{Finanzkennzahlen (\texttt{fin\_*})}

Datei: \file{src/secpd/features/financial.py}.
Rohwerte kommen aus \file{src/secpd/data/edgar.py} (\texttt{TAG\_MAP}, us-gaap).
Divisionen sind NaN-sicher (kein Inf).

\subsubsection{Level --- klassischer PD-Katalog}

\begin{longtable}{@{}p{0.28\textwidth} p{0.32\textwidth} p{0.32\textwidth}@{}}
\toprule
Feature & Formel / Inhalt & Warum (Sprechen) \\
\midrule
\endhead
\texttt{fin\_leverage} & Verbindlichkeiten / Aktiva & Verschuldungsgrad, Altman/Ohlson-Kern \\
\texttt{fin\_debt\_to\_equity} & (DLTT + DLC) / Equity & Kapitalstruktur; negativ bei Distress \\
\texttt{fin\_current\_ratio} & Umlaufvermögen / kurzfr.\ Verb. & kurzfristige Zahlungsfähigkeit \\
\texttt{fin\_quick\_ratio} & (UV $-$ Inventar) / kfr.\ Verb. & ohne Lager, strenger \\
\texttt{fin\_cash\_ratio} & Cash / kfr.\ Verb. & härteste Liquidität \\
\texttt{fin\_wc\_to\_assets} & Working Capital / Aktiva & Altman~$X_1$ \\
\texttt{fin\_roa} & Jahresüberschuss / Aktiva & Ertragskraft \\
\texttt{fin\_net\_margin} & NI / Umsatz & Marge, Preisspielraum \\
\texttt{fin\_interest\_coverage} & EBIT / Zinsaufwand & Zins tragbar? Coverage-Stress \\
\texttt{fin\_re\_to\_assets} & Gewinnrücklagen / Aktiva & Altman~$X_2$, kumulierte Substanz \\
\texttt{fin\_asset\_turnover} & Umsatz / Aktiva & Effizienz, Altman~$X_5$ \\
\texttt{fin\_log\_assets} & $\log(1+\text{Aktiva})$ & Größe / Skaleneffekt \\
\texttt{fin\_d\_log\_revenue} & $\Delta\log$ Umsatz YoY & Wachstumsschock \\
\texttt{fin\_altman\_z} & $1{,}2X_1+1{,}4X_2+3{,}3X_3+0{,}6X_4+X_5$ & Komposit; $X_4$ mit \emph{Buch}-Equity \\
\bottomrule
\end{longtable}

\begin{sprech}
Altman-Z hier mit Buch- statt Marktequity --- wir haben keine Marktpreise.
Trotzdem nützlich als verdichteter Distress-Score. Einzelratios bleiben im
Modell, der RF darf beides nutzen.
\end{sprech}

\subsubsection{Dynamik --- Verfall ist informativer als Niveau}

Für die Basisgrößen Leverage, ROA, D/E, Current Ratio, $\log$ Assets, Turnover,
Net Margin:

\begin{itemize}[leftmargin=1.2em]
  \item \texttt{fin\_d\_*} --- Jahresveränderung (erste Differenz, je CIK).
  \item \texttt{fin\_vol3\_*} --- rollierende 3-Jahres-Standardabweichung
        (min.\ 2 Beobachtungen).
\end{itemize}

\textbf{Warum:} Eine Current Ratio von 1{,}1 kann normal sein. Der \emph{Absturz}
von 2{,}0 auf 1{,}1 in einem Jahr ist das PD-Signal. Volatilität fängt
erratisches Reporting / Geschäftsrisiko.

\subsubsection{Branche --- SIC-2 $\times$ Jahr}

\texttt{fin\_leverage\_sic\_pct}, \texttt{fin\_roa\_sic\_pct},
\texttt{fin\_current\_ratio\_sic\_pct}, \texttt{fin\_wc\_to\_assets\_sic\_pct}:
Perzentilrang in der Zelle SIC-2-Steller $\times$ Geschäftsjahr, nur wenn
$\ge 5$ Beobachtungen, sonst NaN.

\textbf{Warum:} 2$\times$ Leverage in Utilities ist etwas anderes als in Software.
Ohne Peer-Normierung lernt der RF Industrie-Dummies über den Umweg der Niveaus.

\subsubsection{Missing-Indikatoren}

Zu jeder Kennzahl \texttt{fin\_miss\_<name>} $\in\{0,1\}$.

\textbf{Warum:} Die Pipeline imputiert den Median. Ohne Flag wäre ``kein Cash
in XBRL'' unsichtbar --- oft selbst ein Qualitätssignal (junge Filings, Lücken,
nicht-standard Tags).

In den Logs: $\approx 64$ Finanz-Spalten, davon $\approx 13$ Level.

% -----------------------------------------------------------------------------
\subsection{8-K-Event-Features (\texttt{evt\_*})}

Datei: \file{src/secpd/data/events.py}, Funktion \texttt{add\_event\_features}.
Fenster: $(\text{10-K-Filing}-365\,\text{Tage},\;\text{10-K-Filing}]$.
Kein Event $\Rightarrow 0$ (Abwesenheit ist Information, nicht Missing).

\begin{tabularx}{\textwidth}{@{}l l X@{}}
\toprule
Feature & 8-K-Item & Warum \\
\midrule
\texttt{evt\_n\_8k} & alle 8-K/8-K/A & Unruhe, Kommunikationsfrequenz \\
\texttt{evt\_n\_auditor\_change} & 4.01 (alt: 4) & Prüferwechsel oft vor Stress \\
\texttt{evt\_n\_officer\_departure} & 5.02 (alt: 6) & CFO/CEO-Abgang \\
\texttt{evt\_n\_covenant\_accel} & 2.04 & Covenant-Bruch / Beschleunigung --- direkt PD-nah \\
\texttt{evt\_n\_impairment} & 2.06 & Wertberichtigung, Asset-Qualität \\
\texttt{evt\_n\_delisting} & 3.01 & Börsenhinweis; \emph{raus}, wenn Delisting Label ist \\
\bottomrule
\end{tabularx}

\textbf{Nicht im Feature-Raum:}
\begin{itemize}[leftmargin=1.2em]
  \item Insolvenz Item 1.03 --- Zielvariable.
  \item Restatement 4.02 --- Default aus, zu nah an AAER/Fraud.
\end{itemize}

\begin{sprech}
``PIT-sauber'' heißt: nur Informationen, die am 10-K-Tag schon öffentlich waren.
Ein 8-K drei Tage nach dem Filing darf den Score dieses Firm-Years nicht
beeinflussen.
\end{sprech}

% -----------------------------------------------------------------------------
\subsection{Text-Features (nur Combined)}

Datei: \file{src/secpd/features/textual.py}.
Design: \textbf{precompute-then-join} --- kein LLM im sklearn-Transformer.
Gründe: Scoring ohne Netz, Batch/Cache an einer Stelle, kein Fitting also kein Leak.

\subsubsection{LLM --- Vertrag \texttt{TextRiskProfile}}

Datei: \file{src/secpd/llm/schema.py}. Jeder Client (Mock, OpenAI, Bank) muss
dieses Pydantic-Objekt liefern.

Das LLM \emph{extrahiert} mehr, als das Modell \emph{nutzt}:

\begin{tabularx}{\textwidth}{@{}l l X@{}}
\toprule
Feld & Im RF? & Rolle \\
\midrule
\texttt{llm\_risk\_sentiment} $\in[-1,1]$ & \textbf{ja} & Krisenton vs.\ Zuversicht; einziges LLM-Modellfeature \\
\texttt{llm\_vagueness\_score} & nein & Hedging; empirisch ohne Extra-Signal \\
\texttt{llm\_redundancy\_score} & nein & Boilerplate \\
\texttt{llm\_complexity\_score} & nein & früher im Modell, später gestrichen \\
\texttt{llm\_confidence} & nein & Fallback hat 0 --- Debug, nicht Target \\
\texttt{obfuscation\_indicators} & nein & Stichworte, nur Cache/UX \\
\texttt{risk\_summary} & nein & max.\ 3 Sätze, UI \\
\bottomrule
\end{tabularx}

Input: MD\&A, Tabellen raus, Head+Tail auf 12k Zeichen (Anfang = Lage, Ende =
Liquidität/Going-Concern). Cache-Schlüssel = SHA-256 des vorbereiteten Texts.

\begin{sprech}
Warum nur Sentiment? Vagheit und Komplexität waren im Luna-Cache, haben im
Combined-RF aber nichts über Keywords + Sentiment hinaus geliefert. Schlank
halten reduziert Overfit bei wenigen Defaults.
\end{sprech}

\subsubsection{Keyword-Zähler auf dem \emph{vollen} MD\&A}

Nicht auf dem 12k-Ausschnitt --- Going-Concern steht oft hinten.

\begin{tabularx}{\textwidth}{@{}l X@{}}
\toprule
Feature & Typische Phrasen (Zähler, nicht 0/1) \\
\midrule
\texttt{txt\_going\_concern} & going concern, substantial doubt, continue as a going concern \\
\texttt{txt\_liquidity\_stress} & insufficient liquidity, working capital deficit/shortfall, unable to meet \ldots debt \\
\texttt{txt\_covenant} & covenant violation/breach/waiver/default \\
\texttt{txt\_restructuring} & restructur(e|ing), recapitalization, debt exchange/forbearance \\
\texttt{txt\_bankruptcy\_lang} & chapter 11, bankruptcy, receivership \\
\bottomrule
\end{tabularx}

\textbf{Warum Zähler:} Der RF kann Schwellen lernen (ein ``restructuring'' im
Boilerplate vs.\ zehn im Liquiditätsabschnitt).

\begin{sprech}
Im erweiterten Universum kam der Text-Mehrwert vor allem von diesen Keywords,
nicht vom LLM --- der Luna-Cache deckt die neuen CIKs noch dünn. Combined
schlägt Financial trotzdem, weil Going-Concern und Covenant harte, seltene
Flags sind, die Ratios nicht ersetzen.
\end{sprech}

% =============================================================================
\section{Wie das Modell entscheidet}
% =============================================================================

\subsection{Split und Evaluation}

\begin{itemize}[leftmargin=1.2em]
  \item \textbf{Temporal} (Default): letzte Geschäftsjahre = Test. Entspricht Deployment.
  \item \textbf{Group-Split} über CIK: dieselbe Firma nie in Train und Test
        (Memorization-Check).
  \item \textbf{Rolling-Origin}: Cutoffs 2012, 2014, \ldots je 2-Jahres-Testfenster ---
        ehrlicher als ein Split, weil Firm-Overlap im Temporal-Split hoch ist
        ($\approx 80\,\%$ der Testfirmen waren schon im Training).
\end{itemize}

Metriken in \file{src/secpd/evaluation.py}:

\begin{tabularx}{\textwidth}{@{}l X@{}}
\toprule
Metrik & Wann du sie nennst \\
\midrule
ROC-AUC & Ranking-Güte; 0{,}5 Zufall, $>$0{,}7 stark. Bei 1\,\% Basisrate zu schmeichelhaft. \\
PR-AUC & \textbf{Leitmetrik} bei seltenen Events. \\
Top-10\,\%-Capture / Lift & Kreditrisiko-Sprache: ``von 10 Defaults liegen $x$ im riskantesten Dezil''. \\
Brier + Skill vs.\ Basisrate & Kalibrierung; Skill $>0$ = besser als konstante PD. \\
CIK-Cluster-Bootstrap & Konfidenz, weil mehrere Jahre derselben Firma korrelieren. \\
\bottomrule
\end{tabularx}

\subsection{Was du zu den Zahlen sagen kannst}

\textbf{Kleines Label-Set} (533 CIKs, Policy-Gewinner h36 Bankruptcy):
Combined ROC $\approx 0{,}90$, PR $\approx 0{,}30$, Top-10\,\%-Capture $\approx 60\,\%$,
Brier-Skill $+0{,}07$. Group-Split hält (ROC $\approx 0{,}85$) --- kein reines Merken.

\textbf{Volles Universum} ($\approx 5\,500$ CIKs, Rolling):
h36 Bankruptcy pooled $\approx 316$ Positive, Combined ROC $\approx 0{,}80$,
Top-10\,\% $\approx 53\,\%$. Delisting-h24 gewinnt mean-PR nur über die Basisrate
($\approx 2\,700$ Positive), Ranking ist schlechter (ROC $\approx 0{,}75$, Capture $\approx 39\,\%$).

\begin{sprech}
Satz zum Delisting: ``Mehr Events $\neq$ besseres Label. PR-AUC steigt mit der
Basisrate. Für Ranking und Watchlist bleibt 36-Monats-Insolvenz der Gewinner.''
\end{sprech}

% =============================================================================
\section{Code-Karte: wo steht was?}
% =============================================================================

\subsection{Einstiege (was du live zeigen kannst)}

\begin{tabularx}{\textwidth}{@{}l X@{}}
\toprule
Datei & Rolle \\
\midrule
\file{start.py} & Interaktives Menü: scoren, Güte, LLM-Key, Fetch, Training \\
\file{train.py} & Kanonischer Trainingspfad (Bank und Home identisch) \\
\file{predict.py} & Minimales Serving: Bundle laden, PD schreiben \\
\file{README.md} & Architektur, ENV, Home$\rightarrow$Bank \\
\file{requirements.txt} & Pins; Home und Bank \emph{identisch} \\
\bottomrule
\end{tabularx}

\subsection{Paket \texttt{src/secpd/}}

\begin{tabularx}{\textwidth}{@{}l X@{}}
\toprule
Datei & Was du dort findest \\
\midrule
\file{config.py} & ENV: \texttt{SECPD\_LLM\_*}, \texttt{SECPD\_SEC\_UA}, Cache-Pfad \\
\file{splitting.py} & \texttt{smart\_split}, \texttt{rolling\_origin\_splits} \\
\file{evaluation.py} & ROC/PR/Brier, Dezile, Cluster-Bootstrap \\
\file{data/zenodo.py} & JSON streamen, AAER-Join, Spalten auflösen \\
\file{data/edgar.py} & \texttt{TAG\_MAP}, companyfacts $\rightarrow$ Panel \\
\file{data/events.py} & Item-Map, Label, PIT-Features, Leakage-Guard \\
\file{data/synthetic.py} & Demo-Daten mit echtem Signalfluss (ohne Download) \\
\file{features/financial.py} & alle \texttt{fin\_*} \\
\file{features/textual.py} & \texttt{prepare\_text}, Keywords, LLM-Join \\
\file{llm/schema.py} & \texttt{TextRiskProfile} --- der Vertrag \\
\file{llm/base.py} & ABC + Batch-Fallback \\
\file{llm/mock.py} & deterministische Heuristik, Tests/Demo \\
\file{llm/bank.py} & OpenAI-kompatibles \texttt{/v1/chat/completions} \\
\file{llm/cache.py} & JSON je Text-Hash, \texttt{cache\_only} \\
\file{llm/\_\_init\_\_.py} & Factory \texttt{get\_llm\_client()} \\
\file{models/pipeline.py} & RF + Imputer + Kalibrierung; Text-LogReg \\
\file{models/ensemble.py} & Logit-Mittelung Option~B \\
\file{models/persistence.py} & joblib + Versionsstempel \\
\bottomrule
\end{tabularx}

\subsection{Skripte}

\begin{tabularx}{\textwidth}{@{}l X@{}}
\toprule
Skript & Wann \\
\midrule
\file{scripts/fetch\_zenodo.py} & Roh-MD\&A / AAER laden \\
\file{scripts/convert\_zenodo.py} & $\rightarrow$ \file{zenodo\_labeled.csv.gz} bzw.\ \file{zenodo\_full.csv.gz} \\
\file{scripts/fetch\_edgar\_financials.py} & companyfacts-Panel, resumefähiger JSON-Cache \\
\file{scripts/fetch\_edgar\_events.py} & 8-K-Metadaten \\
\file{scripts/precompute\_llm\_features.py} & Luna-Cache füllen (wiederaufnehmbar) \\
\file{scripts/rolling\_eval.py} & Financial vs.\ Combined über Cutoffs \\
\file{scripts/compare\_policies.py} & $H\in\{12,24,36\}\times$ bankruptcy[/delisting] \\
\file{scripts/freeze\_benchmark.py} & reproduzierbarer Frozen-Report + Hash \\
\file{scripts/build\_offline\_bundle.sh} & manylinux-Wheels für die Bank (3.14: \texttt{PLATFORM=manylinux\_2\_28\_x86\_64}) \\
\file{scripts/install\_offline.sh} & pip \texttt{--no-index} auf der Box \\
\bottomrule
\end{tabularx}

\subsection{Tests --- was du als Qualitätssicherung erwähnst}

\begin{tabularx}{\textwidth}{@{}l X@{}}
\toprule
Test & Absicherung \\
\midrule
\file{tests/test\_events.py} & Label-Fenster, Leakage (1.03 nicht in Features), End-to-End Default \\
\file{tests/test\_financial\_features.py} & Ratios, NaN-sichere Division \\
\file{tests/test\_keyword\_features.py} & Distress-Zähler \\
\file{tests/test\_splitting\_rolling.py} & Temporal / Rolling ohne Klassenkollaps \\
\file{tests/test\_evaluation.py} & Metriken, Dezile \\
\file{tests/test\_llm\_mock.py} & Schema-Vertrag, Mock deterministisch \\
\file{tests/test\_end\_to\_end.py} & Synth $\rightarrow$ Fit $\rightarrow$ joblib \\
\bottomrule
\end{tabularx}

33 Tests, ohne Netz.

\subsection{Daten- und Artefaktpfade}

\begin{lstlisting}
data/raw/firm_years.json              # volles MD&A, lokal, nicht im Git
data/raw/aaer_mark5.csv               # Fraud-Fenster
data/raw/financials_panel[_full].csv  # XBRL-Panel (committbar)
data/raw/edgar_8k_events[_full].csv   # 8-K-Metadaten
data/raw/edgar_companyfacts/          # JSON-Cache, lokal
data/raw/edgar_submissions/           # JSON-Cache, lokal
data/processed/zenodo_labeled.csv.gz  # 533 CIKs
data/processed/zenodo_full.csv.gz     # ~5487 CIKs, gitignore (zu groß)
data/cache/llm/bank-gpt-5.6-luna/     # Luna-Replay
models/*_default_h{12,36}.joblib
benchmarks/rolling_* / policy_compare*.md
\end{lstlisting}

% =============================================================================
\section{Live-Demo: welche Datei aufklappen?}
% =============================================================================

Falls jemand den Laptop sieht --- in dieser Reihenfolge:

\begin{enumerate}[leftmargin=1.4em]
  \item \file{src/secpd/data/events.py} --- \texttt{ITEM\_MAP}: ``1.03 nur Label''.
  \item \file{src/secpd/features/financial.py} --- der Ratio-Katalog.
  \item \file{src/secpd/features/textual.py} --- \texttt{\_KEYWORD\_GROUPS}.
  \item \file{src/secpd/llm/schema.py} --- \texttt{MODEL\_FEATURE\_FIELDS = (risk\_sentiment,)}.
  \item \file{src/secpd/models/pipeline.py} --- RF + Kalibrierung.
  \item \file{benchmarks/policy\_compare.md} (klein) und
        \file{benchmarks/policy\_compare\_full.md} (Universum).
\end{enumerate}

Training (Produktion, Cache-only):
\begin{lstlisting}
python train.py \
  --data data/processed/zenodo_full.csv.gz \
  --financials data/raw/financials_panel_full.csv \
  --events data/raw/edgar_8k_events_full.csv \
  --label-source default --default-horizon 36 \
  --label-concepts bankruptcy --require-financials \
  --mode combined --llm openai --llm-cache-only --calibrate
\end{lstlisting}

% =============================================================================
\section{Einwände und klare Antworten}
% =============================================================================

\begin{description}[style=nextline,leftmargin=1.2em,font=\bfseries]
  \item[Ist das eine regulatorische PD?]
    Nein. Es ist ein Distress-/Insolvenz-Proxy-Score aus öffentlichen Filings.
    Kalibriert auf 8-K-Item~1.03, nicht auf IRB-Default im Kreditvertragssinn.
  \item[Warum nicht Gradient Boosting?]
    Auf dem Gewinner-Policy hat RF mit Sigmoid-Kalibrierung GBM geschlagen
    (wenige Positive, $\approx 100$ im kleinen Set). Hebel ist Datenuniversum,
    nicht der Learner.
  \item[Look-ahead / Firm-Leakage?]
    Temporal + Rolling-Origin; Group-Split als Holdout; Eventfenster PIT;
    Insolvenz-Item ausgeschlossen. Tests dazu.
  \item[Warum MD\&A und nicht der ganze 10-K?]
    MD\&A ist der narrative Risikoteil. Tabellen im Zenodo-Dump sind Zahlensuppe
    und werden vor dem LLM gestrichen.
  \item[Untererfassung von Defaults?]
    Ja. Nicht jede Insolvenz erzeugt ein 8-K (Delisting vorher, Filing-Pflichten
    weg). Deshalb PR-AUC, nicht nur ROC; und deshalb Delisting als Label getestet
    und verworfen.
  \item[Python 3.14 auf der Bank?]
    Alte Pins ohne cp314-Wheels. Aktuell: numpy~2.3.5, pandas~2.3.3, sklearn~1.7.2.
    Modelle danach neu fitten.
  \item[Kann das interne LLM rein?]
    \texttt{get\_llm\_client()} / \file{llm/bank.py}: Endpoint, Key, Modell über ENV.
    Marker \texttt{\# ADAPT}. Cache macht Home und Bank bitgleich.
\end{description}

% =============================================================================
\section{Ein-Folien-Gerüst (wenn du slides baust)}
% =============================================================================

\begin{enumerate}[leftmargin=1.4em]
  \item Problem: öffentliche Filings $\rightarrow$ früher Distress-Score, air-gapped.
  \item Drei Quellen, ein Join \texttt{(cik, fyear)}; 1.03 nur Label.
  \item Features: Ratios+Trends+SIC, 8-K-Frühwarnung, MD\&A-Keywords+Sentiment.
  \item Modell: RF, kalibriert, Combined vs.\ Financial-Baseline.
  \item Güte: Rolling-PR, Top-10\,\%-Capture, Group-Split.
  \item Policy: h36 Bankruptcy, Legacy vor 2004 aus, ab GJ 2009.
  \item Übergabe: Cache, joblib-Stempel, Offline-Wheels, 3.14-Pins.
\end{enumerate}

\vspace{1.2em}
\noindent\textcolor{secpdblue}{\rule{\textwidth}{0.6pt}}\\[0.4em]
{\small
Dieser Zettel beschreibt den Code-Stand auf \texttt{feat/distress-text-features}.
Zahlen aus \file{benchmarks/policy\_compare.md} und
\file{benchmarks/policy\_compare\_full.md} bei der Präsentation aufschlagen,
nicht aus dem Kopf zitieren, wenn die Reports sich geändert haben.}

\end{document}
